\documentclass[11pt, letterpaper]{article}
\usepackage{mobi-lectura}

\title{Diseño de Experimentos de Simulación con Monte Carlo}
\author{Alejandra Tabares Pozos}
\date{}

\begin{document}

\maketitle
\tableofcontents
\newpage

\section{Introducción: del Fundamento al Diseño}

La lectura anterior (El Teorema del Límite Central) estableció el fundamento teórico: la media de $n$ réplicas independientes $\bar{X}_n$ converge en distribución a una Normal, lo que permite construir un intervalo de confianza de la forma:
\begin{equation}
\bar{X}_n \;\pm\; t_{n-1,\,1-\alpha/2}\,\frac{S}{\sqrt{n}}
\end{equation}
Esta lectura se enfoca en la \textbf{dimensión práctica}: ¿cómo diseñar el experimento de simulación para obtener la precisión requerida? Las preguntas concretas son: ¿cuántas réplicas son suficientes?, ¿cómo sé si mis resultados son lo bastante precisos para tomar una decisión?, y ¿cómo comunico los resultados con rigor estadístico?

\section{Fórmula de diseño: número de réplicas}

El ancho del IC depende de $n$: más réplicas $\Rightarrow$ SE más pequeño $\Rightarrow$ IC más angosto. La relación es:
\[
\text{Margen de error} = t_{n-1,\,1-\alpha/2}\,\frac{S}{\sqrt{n}} \approx z_{1-\alpha/2}\,\frac{\sigma}{\sqrt{n}}
\]
Para garantizar un margen máximo $E$ con confianza $1-\alpha$, despejamos $n$:
\begin{equation}
n \ge \left(\frac{z_{1-\alpha/2}\,\sigma}{E}\right)^2
\end{equation}
En la práctica $\sigma$ es desconocida: se usa una estimación piloto $s$ de una corrida inicial de 10--20 réplicas, y se redondea $n$ hacia arriba.

\begin{metodobox}[Protocolo de diseño en 4 pasos]
\begin{enumerate}
    \item \textbf{Corrida piloto}: ejecutar $n_0 = 10\text{--}20$ réplicas; calcular $\bar{X}_{n_0}$ y $S_{n_0}$.
    \item \textbf{Determinar $n^*$}: aplicar la fórmula con $\sigma \approx S_{n_0}$ y el margen de error deseado $E$.
    \item \textbf{Completar réplicas}: si $n^* > n_0$, ejecutar $n^* - n_0$ réplicas adicionales.
    \item \textbf{Reportar IC}: calcular el IC final con todas las $n^*$ réplicas.
\end{enumerate}
\end{metodobox}

\section{Ejemplo 1: Introductorio - Estimación de Costo de Materia Prima}
Un analista financiero desea estimar el costo promedio de una orden de compra. El precio y la cantidad son aleatorios. Se realizan 50 simulaciones independientes y se obtienen los siguientes estadísticos:
\begin{itemize}
    \item Media muestral ($\bar{X}$): \$1,250 USD
    \item Desviación estándar muestral ($S$): \$200 USD
    \item Tamaño de muestra ($n$): 50
\end{itemize}

Para un nivel de confianza del 95\% ($\alpha = 0.05$):
1. Grados de libertad: $50 - 1 = 49$.
2. Valor crítico $t_{49, 0.975} \approx 2.009$.
3. Error estándar: $200 / \sqrt{50} \approx 28.28$.
4. Margen de error: $2.009 \times 28.28 \approx 56.82$.

\textbf{Resultado:} El costo promedio esperado está entre \textbf{\$1,193.18 y \$1,306.82} con un 95\% de confianza.

\section{Ejemplo 2: Nivel Medio - Análisis de Riesgo en Proyectos}
El costo total de un proyecto es $CT = L + M + C$, donde las variables son:
\begin{itemize}
    \item Mano de obra: $L \sim \text{Triangular}(a=10, b=20, \text{moda}=15)$
    \item Materiales: $M \sim N(\mu=50, \sigma=5)$
    \item Contingencias: $C \sim U(5, 15)$
\end{itemize}
Por la linealidad del valor esperado y la independencia de las variables:
\[ E[CT] = E[L] + E[M] + E[C] = \frac{10+15+20}{3} + 50 + \frac{5+15}{2} = 15 + 50 + 10 = 75 \text{ (unidades)} \]
La varianza total (por independencia) es:
\[ \text{Var}[CT] = \text{Var}[L] + \text{Var}[M] + \text{Var}[C] = \frac{10^2+15^2+20^2-10\cdot15-15\cdot20-10\cdot20}{18} + 25 + \frac{100}{12} \]
\[ = \frac{100+225+400-150-300-200}{18} + 25 + 8.33 = \frac{75}{18} + 33.33 \approx 4.17 + 33.33 = 37.5 \]

Tras $n = 10{,}000$ réplicas, el error estándar del estimador $\bar{CT}$ sería:
\[ SE = \frac{\sigma_{CT}}{\sqrt{n}} = \frac{\sqrt{37.5}}{\sqrt{10000}} = \frac{6.12}{100} = 0.0612 \text{ unidades} \]
El IC del 95\% resultante tendría un margen de error de $\pm 1.96 \times 0.0612 \approx \pm 0.12$, es decir, el costo promedio estimado sería $75 \pm 0.12$ con alta precisión. Para reducir el error a la mitad bastarían $n = 2{,}500$ réplicas menos (dado que $SE \propto 1/\sqrt{n}$, duplicar la precisión requiere cuadruplicar el número de réplicas).


\section{Interpretación y Errores Comunes}
Es vital entender que el intervalo de confianza \textbf{no} contiene el 95\% de los datos de los costos (ese sería un intervalo de predicción). El IC del 95\% significa que, si repitiéramos este experimento de simulación 100 veces, en 95 de esas veces el intervalo calculado contendría el verdadero valor medio del sistema.

Si el ancho del intervalo es demasiado grande para la toma de decisiones, el ingeniero tiene una sola vía: \textbf{aumentar el número de réplicas ($n$)}. Dado que el error disminuye con $\sqrt{n}$, para reducir el error a la mitad se requiere cuadruplicar el número de simulaciones.

\section*{Referencias}
\begin{enumerate}[label={[\arabic*]}]
    \item Law, A.\,M. (2015). \textit{Simulation Modeling and Analysis}, 5th ed. McGraw-Hill. Cap.~9.
    \item DeGroot, M.\,H. \& Schervish, M.\,J. (2012). \textit{Probability and Statistics}, 4th ed. Pearson. Cap.~8.
    \item Banks, J., Carson, J., Nelson, B., \& Nicol, D. (2014). \textit{Discrete-Event System Simulation}, 5th ed. Pearson. Cap.~10.
    \item Casella, G. \& Berger, R.\,L. (2002). \textit{Statistical Inference}, 2nd ed. Duxbury. Cap.~9.
\end{enumerate}

\end{document}
